\begin{tikzpicture}[
  font=\sffamily,
  atom/.style={
    rectangle, rounded corners=3pt, draw=ceindigo, line width=0.8pt,
    fill=cepanel, text width=3.1cm, align=left, inner sep=5pt, minimum height=1.75cm, font=\sffamily\scriptsize
  }
]
  \node[font=\sffamily\bfseries\small, ceindigo] at (5.3, 3.35)
    {LOS OCHO MODELOS DE NEGOCIO ELECTRÓNICOS ATÓMICOS (WEILL Y VITALE)};

  % Fila 1
  \node[atom] (w1) at (0, 1.8)
    {\textcolor{ceindigo}{\bfseries 1. Directo al cliente}\\Evita intermediarios; controla relación y datos (Dell, Zara web).};
  \node[atom] (w2) at (3.55, 1.8)
    {\textcolor{ceindigo}{\bfseries 2. Proveedor de contenidos}\\Monetiza información, cursos o media digital (Spotify, prensa).};
  \node[atom] (w3) at (7.1, 1.8)
    {\textcolor{ceindigo}{\bfseries 3. Intermediario}\\Mercado o agregador por comisión de transacción (Booking, Wallapop).};
  \node[atom] (w4) at (10.65, 1.8)
    {\textcolor{ceindigo}{\bfseries 4. Servicio completo}\\Un único punto resuelve necesidad financiera o vital integral.};

  % Fila 2
  \node[atom, draw=ceteal, fill=cepanelteal] (w5) at (0, -0.45)
    {\textcolor{ceteal!90!black}{\bfseries 5. Infraestructura compartida}\\Consorcio de competidores que comparten red común (SWIFT, Bizum).};
  \node[atom, draw=ceteal, fill=cepanelteal] (w6) at (3.55, -0.45)
    {\textcolor{ceteal!90!black}{\bfseries 6. Integrador red de valor}\\Coordina datos y logística de múltiples socios sin fabricar (Nike, Uber).};
  \node[atom, draw=ceteal, fill=cepanelteal] (w7) at (7.1, -0.45)
    {\textcolor{ceteal!90!black}{\bfseries 7. Comunidad virtual}\\Miembros fidelizados alrededor de un interés común (GitHub, Strava).};
  \node[atom, draw=ceteal, fill=cepanelteal] (w8) at (10.65, -0.45)
    {\textcolor{ceteal!90!black}{\bfseries 8. Entidad pública / B2G}\\Sede electrónica y trámites abiertos con alta confianza jurídica.};
\end{tikzpicture}